\chapter{Complexity and Constraints}\label{ch:complexity}

A correct solution that is too slow scores exactly the same as no solution at all. Complexity analysis tells you, before
you write a line, whether an idea can pass; it is the tool you use in step 2 of the problem-solving loop to throw away
hopeless ideas and to recognise which technique the problem setter had in mind. This chapter teaches you to estimate the
running time and memory of an algorithm quickly and reliably, and it gives you real measured costs of the operations you
will use every day, so that ``fast enough'' stops being a guess.

\begin{chapterintro}{\bronze\ \silver\ \gold}{Chapter~\ref{ch:approach}.}
  \item count operations and express them in $\mathcal{O}$-notation;
  \item read a problem's limits and deduce the complexity it expects;
  \item analyse loops, including harmonic sums, square-root loops and amortized costs;
  \item solve recurrences with recursion trees and the master theorem;
  \item budget memory, and know what each C++ operation really costs, measured.
\end{chapterintro}

% ======================================================================================================================
\section[Counting operations]{Counting operations \hfill \bronze}\label{sec:counting-ops}

A modern judge executes on the order of $10^8$ to $10^9$ \emph{simple operations} per second: additions, comparisons,
array accesses, branches. The exact number depends on the machine and on what the operations are (Section~\ref{sec:measured}
measures them), but for planning purposes one rule is enough:

\begin{fact}[The budget]
Plan for about $10^8$ simple operations per second of time limit. Below $10^7$ you are safe by a wide margin; around
$10^8$ you are fine if the inner loop is simple; above $10^9$ you will not pass without a very special trick.
\end{fact}

So the question ``will it pass?'' becomes ``how many operations does it do on the largest input?''. For loops that is
mostly multiplication. The code below does work once for every pair $i < j$:

\begin{code}
for (int i = 0; i < n; i++)          // n iterations
    for (int j = i + 1; j < n; j++)  // n-1, n-2, ..., 1, 0 iterations
        work(i, j);                  // in total n(n-1)/2 calls
\end{code}

For $n = 5000$ that is $1.25 \cdot 10^7$ calls: instant. For $n = 2\cdot10^5$ it is $2 \cdot 10^{10}$: over a minute on any
judge. The same idea is fine or hopeless depending on a single number in the statement, which is why you always compute
the count \emph{with the real limit plugged in}.

Counting is rarely exact, and it does not need to be: you only need the order of magnitude. Three habits make it quick.
First, count the innermost operation and multiply by the number of times each enclosing loop runs. Second, round
generously: $\log_2(2\cdot10^5) \approx 18$, $\sqrt{2\cdot10^5} \approx 450$, $2^{20} \approx 10^6$. Third, remember that a
``simple operation'' inside a library call may itself be a loop: \cc{s.count(x)} on a \cc{vector} is $n$ operations, not one.

% ======================================================================================================================
\section[Big-O notation]{Big-O notation \hfill \bronze}\label{sec:bigo}

Counting exactly is tedious and machine-dependent, so we describe how the count \emph{grows} with the input size and ignore
the rest. We say that an algorithm runs in time \defn{$\mathcal{O}(f(n))$} if, for all large enough $n$, it performs at
most $c \cdot f(n)$ operations for some fixed constant $c$. Two companions complete the vocabulary: $\Omega(f(n))$ means ``at
least $c\cdot f(n)$'', and $\Theta(f(n))$ means both, i.e.\ exactly this order of growth. In contests people say
$\mathcal{O}$ even when they mean $\Theta$; the meaning is usually clear.

The rules follow from the definition. Constant factors disappear: $5n$ and $n / 3$ are both $\Oh{n}$. Lower-order terms
disappear: $3n^2 + 10n + 7$ is $\Oh{n^2}$, because for large $n$ the $n^2$ term dominates everything else. Logarithm bases
disappear too, since $\log_2 n = \log_2 10 \cdot \log_{10} n$ differs by a constant factor. With several inputs you keep
each of them: BFS on a graph is $\Oh{n + m}$, and neither term may be dropped, because either can be the larger one.

\begin{figure}[htbp]
\centering
\begin{tikzpicture}
\begin{loglogaxis}[cpplot, width=11.8cm, height=7.2cm, xmin=1, xmax=1e7, ymin=1, ymax=1e13, xlabel={input size $n$},
    ylabel={operations}, legend pos=outer north east, samples=120]
  % parametric in t = log10(n), so that the samples are spread evenly on the logarithmic axis
  \addplot[fgverm, dashed, line width=0.9pt, domain=0:7, forget plot] ({10^x}, {1e8});
  \addplot+[mark=none, domain=0.302:7] ({10^x}, {x*ln(10)/ln(2)}); \addlegendentry{$\log n$}
  \addplot+[mark=none, domain=0:7] ({10^x}, {10^x}); \addlegendentry{$n$}
  \addplot+[mark=none, domain=0.302:7] ({10^x}, {10^x*x*ln(10)/ln(2)}); \addlegendentry{$n \log n$}
  \addplot+[mark=none, domain=0:6.5] ({10^x}, {10^(2*x)}); \addlegendentry{$n^2$}
  \addplot+[mark=none, domain=0:4.33] ({10^x}, {10^(3*x)}); \addlegendentry{$n^3$}
  \addplot+[mark=none, domain=0:1.63] ({10^x}, {2^(10^x)}); \addlegendentry{$2^n$}
  \addplot+[mark=none] coordinates {(1,1) (2,2) (3,6) (4,24) (5,120) (6,720) (7,5040) (8,40320) (9,362880)
    (10,3628800) (11,39916800) (12,479001600) (13,6227020800) (14,87178291200) (15,1307674368000)}; \addlegendentry{$n!$}
  \node[font=\sffamily\scriptsize\bfseries, text=fgverm, anchor=south east] at (axis cs:1e7, 1e8) {$10^8$: one second};
\end{loglogaxis}
\end{tikzpicture}
\caption{How fast common complexities grow (both axes logarithmic). Each curve crosses the one-second budget (dashed) at the
largest $n$ it can handle: $n!$ near 11, $2^n$ near 27, $n^3$ near 460, $n^2$ near $10^4$, $n\log n$ near $4.5\cdot10^6$.
Polynomials are straight lines in this view; exponentials bend upwards and escape almost immediately.}
\label{fig:growth}
\end{figure}

\begin{pitfall}[Big-O hides constants, and constants can be big]
$\Oh{n}$ with a constant of 100 is slower than $\Oh{n \log n}$ with a constant of 1 for every $n$ you will meet. Big-O tells
you which of two ideas \emph{scales} better, not which is faster on today's input. When two ideas have the same complexity,
the constant factor (Section~\ref{sec:measured}) decides.
\end{pitfall}

\paragraph{Worst case, average case, amortized.} Unless stated otherwise, complexity means the \emph{worst case}: the most
operations over all inputs of size $n$. That is what matters in contests, because the tests are chosen by someone trying to
break you. \emph{Expected} (average) complexity appears with randomised algorithms: quicksort with a random pivot, hash
tables, treaps. \emph{Amortized} complexity bounds the total cost of a sequence of operations even when single operations
are occasionally expensive (Section~\ref{sec:amortized}).

% ======================================================================================================================
\section[What the limits tell you]{What the limits tell you \hfill \bronze}\label{sec:constraints}

Problem setters choose limits to allow the intended solution and to kill slower ones. Reading their intention from the
limits is the single most useful skill in this chapter: find the largest input parameter in the statement, look it up in
Table~\ref{tab:constraints}, and you know the target complexity and a short list of techniques before you have had a
single idea. Figure~\ref{fig:maxn} shows where each complexity runs out of budget.

\begin{table}[htbp]
\centering
\caption{Limits and the complexities that fit in about one second.}\label{tab:constraints}
\small
\begin{tabularx}{\linewidth}{>{\bfseries}p{2.3cm}p{4.3cm}L}
\toprule
\thd{Largest $n$} & \thd{Complexity} & \thd{Typical techniques} \\
\midrule
$n \le 10$ & $\Oh{n!}$, $\Oh{n^6}$ & all permutations, brute force over everything \\
$n \le 20$ & $\Oh{2^n \cdot n}$ & all subsets, bitmask DP (Chapter~\ref{ch:dppatterns}) \\
$n \le 40$ & $\Oh{2^{n/2} \cdot n}$ & meet in the middle (Chapter~\ref{ch:search}) \\
$n \le 100$ & $\Oh{n^4}$ & DP with four nested indices \\
$n \le 500$ & $\Oh{n^3}$ & Floyd--Warshall, interval DP, Gaussian elimination \\
$n \le 5000$ & $\Oh{n^2}$ & all pairs, two-dimensional DP, LCS \\
$n \le 10^5$ & $\Oh{n \sqrt n}$ & square-root decomposition, Mo's algorithm \\
$n \le 2 \cdot 10^5$ & $\Oh{n \log n}$, $\Oh{n \log^2 n}$ & sorting, binary search, heaps, segment trees, Dijkstra \\
$n \le 10^6$ & $\Oh{n}$, or $\Oh{n \log n}$ if simple & linear scans, two pointers, prefix sums, sieves \\
$n \le 10^9$ & $\Oh{\sqrt n}$, $\Oh{\log n}$ & trial division, formulas, binary search \\
$n \le 10^{18}$ & $\Oh{\log n}$, $\Oh{\log^2 n}$, $\Oh{1}$ & fast exponentiation, matrix power, digit DP, formulas \\
\bottomrule
\end{tabularx}
\end{table}

\begin{figure}[htbp]
\centering
\begin{tikzpicture}[x=1.25cm]
  % bar length = log10 of the largest n with f(n) <= 1e8
  \foreach \lab/\v/\txt [count=\k] in {{$n!$}/1.04/{$n \le 11$}, {$2^n\,n$}/1.34/{$n \le 22$}, {$n^4$}/2.0/{$n \le 100$},
      {$n^3$}/2.67/{$n \le 464$}, {$n^2$}/4.0/{$n \le 10^4$}, {$n\sqrt n$}/5.33/{$n \le 2\cdot10^5$},
      {$n \log n$}/6.65/{$n \le 4.5\cdot10^6$}, {$n$}/8.0/{$n \le 10^8$}} {
    \pgfmathsetmacro{\y}{-0.55*\k}
    \node[font=\small, anchor=east] at (-0.1, \y) {\lab};
    \fill[fgblue!35, draw=fgblue, line width=0.5pt] (0, \y-0.19) rectangle (\v, \y+0.19);
    \node[font=\sffamily\scriptsize, anchor=west] at (\v+0.05, \y) {\txt};
  }
  \draw[line width=0.6pt, fgink!70] (0, -0.2) -- (0, -4.7);
  \foreach \e in {0,...,8} {
    \draw[line width=0.4pt, fgink!70] (\e, -4.7) -- (\e, -4.78) node[below, lbl] {$10^{\e}$};
    \draw[line width=0.3pt, fgline!40] (\e, -0.2) -- (\e, -4.7);
  }
  \node[lbl] at (4, -5.45) {largest $n$ that fits in $10^8$ operations (logarithmic scale)};
\end{tikzpicture}
\caption{The largest input each complexity can process with about $10^8$ operations. Moving one row up the table costs
orders of magnitude of input size; this is why the limits reveal the intended complexity so reliably.}
\label{fig:maxn}
\end{figure}

\paragraph{Several parameters.} Statements often have more than one limit: $n$ vertices and $m$ edges, $n$ elements and $q$
queries, a grid of $r \times c$ cells. Multiply out the candidate complexity with all of them: $\Oh{(n + q) \log n}$ with
$n = q = 2 \cdot 10^5$ is about $7 \cdot 10^6$; $\Oh{nq}$ is $4\cdot 10^{10}$. The two limits that look alike often play
different roles: a grid with $r \cdot c \le 10^6$ but $r$ up to $10^6$ is a hint that the solution must not assume a square.

\paragraph{The small limit is the hint.} When one limit is suspiciously small (values up to 100, at most 20 special
vertices, an alphabet of size 3, $k \le 10$), the intended solution almost certainly iterates over it. ``$n \le 2\cdot10^5$,
$k \le 20$'' strongly suggests $\Oh{nk}$ or $\Oh{n \cdot 2^k}$, never something independent of $k$.

\paragraph{Sums over test cases.} ``The sum of $n$ over all test cases does not exceed $2\cdot10^5$'' allows
$\Oh{n\log n}$ \emph{per test}. It does not allow work that is independent of $n$: resetting a global array of size $10^6$,
or a loop to $10^6$, costs $10^6$ \emph{per test case}, and there may be $10^4$ tests. Reset only what you used.

\paragraph{Other time limits.} A 2-second limit doubles the budget; half a second halves it. Limits far above one second
(5 or 10 seconds) usually mean the intended solution is heavy, often $\Oh{n\sqrt n}$ or a large constant, and a clean
implementation of it will pass comfortably. Limits far below usually mean that input and output are large, so fast I/O
matters.

\paragraph{Memory limits.} The usual limit is 256\,MB, sometimes 512\,MB or 1\,GB, and occasionally as low as 16--64\,MB to
force a clever idea. Section~\ref{sec:memory} converts megabytes into numbers of elements.

% ======================================================================================================================
\section[Analysing loops]{Analysing loops \hfill \silver}\label{sec:loops}

Most running times are loops, and a small number of loop shapes cover almost everything.

\paragraph{Independent nested loops multiply.} Two loops to $n$ and $m$ cost $n \cdot m$. Three nested loops to $n$ cost
$n^3$ even if the body is a single addition.

\paragraph{Dependent loops sum.} If the inner loop's length depends on the outer variable, add up the lengths. The pair loop
of Section~\ref{sec:counting-ops} sums $0 + 1 + \dots + (n-1) = n(n-1)/2$, still $\Theta(n^2)$; a triple of nested dependent
loops sums to about $n^3 / 6$. Halving the constant is nice but does not change the class.

\paragraph{Loops that halve are logarithmic.} \cc{while (x > 0) x /= 2;} runs $\lfloor \log_2 x \rfloor + 1$ times. Binary
search, fast exponentiation, Euclid's gcd algorithm and binary lifting all shrink the problem by a constant factor per
step, so they take $\Oh{\log n}$ steps. A loop that multiplies by 2 until it passes $n$ is the same loop read backwards.

\paragraph{Square-root loops.} \cc{for (ll d = 1; d * d <= n; d++)} runs $\sqrt n$ times. It enumerates divisors in pairs
$(d, n/d)$ and tests primality by trial division: $10^6$ steps for $n = 10^{12}$.

\paragraph{The harmonic sum.} Loops over multiples are cheaper than they look:

\begin{code}
for (int d = 1; d <= n; d++)
    for (int x = d; x <= n; x += d)  // n / d iterations
        cnt[x]++;                     // then cnt[x] = number of divisors
\end{code}

The inner loop runs $n / d$ times, so the total is $n\,(1 + \frac12 + \frac13 + \dots + \frac1n) = n \cdot H_n \approx n
\ln n$. For $n = 10^6$ that is $1.4 \cdot 10^7$ steps: this ``for each $d$, visit every multiple'' pattern is how the sieve
of Eratosthenes, divisor counting and many number-theory precomputations run in $\Oh{n \log n}$ or better
(Chapter~\ref{ch:numbertheory}). The reverse, ``for each $x$, try every $d$ and check whether it divides'', costs $n^2$.

\paragraph{Loops over submasks.} Iterating over every submask of every mask of $n$ bits looks like $4^n$ but is $3^n$: each
of the $n$ bits is in one of three states (in the submask, in the mask only, in neither). For $n = 15$ that is $1.4 \cdot
10^7$ instead of $10^9$ (Chapter~\ref{ch:dppatterns}).

\begin{tip}
When a loop's cost is not obvious, instrument it: count the iterations in a global variable on the largest input you can
generate and print the count. Measuring beats guessing whenever the analysis gets subtle.
\end{tip}

% ======================================================================================================================
\section[Amortized analysis]{Amortized analysis \hfill \gold}\label{sec:amortized}

Sometimes a single step is expensive but the steps \emph{together} are cheap. Amortized analysis bounds the total cost of a
whole sequence of operations instead of multiplying the worst single step by the number of steps.

\subsection{The aggregate method}

Count the total work directly. The standard example is \cc{vector::push_back}. When the vector is full, it allocates twice
the capacity and copies every element, which costs as much as the current size. Most pushes cost 1, but pushes number 2, 3,
5, 9, 17, \ldots\ copy $1, 2, 4, 8, 16, \ldots$ elements (Figure~\ref{fig:doubling}). Over $n$ pushes the copies total at
most $1 + 2 + 4 + \dots < 2n$, so all $n$ pushes cost less than $3n$: $\Oh{1}$ amortized per push.

\begin{figure}[htbp]
\centering
\begin{tikzpicture}[x=0.32cm, y=0.24cm]
  \foreach \k in {1,...,32} {
    \pgfmathtruncatemacro{\km}{\k-1}
    \pgfmathsetmacro{\cst}{(\k > 1 && (\km == 1 || \km == 2 || \km == 4 || \km == 8 || \km == 16)) ? \k : 1}
    \pgfmathsetmacro{\isbig}{\cst > 1 ? 1 : 0}
    \ifdim\isbig pt>0.5pt \fill[fgverm!45, draw=fgverm, line width=0.4pt] (\k-0.9, 0) rectangle (\k-0.1, \cst);
    \else \fill[fgblue!30, draw=fgblue, line width=0.4pt] (\k-0.9, 0) rectangle (\k-0.1, \cst); \fi
  }
  % running average: total cost so far / number of pushes (computed exactly; the total after 32 pushes is 63)
  \draw[line width=1.3pt, fggreen] plot coordinates {(0.5,1.000) (1.5,1.500) (2.5,2.000) (3.5,1.750) (4.5,2.400)
    (5.5,2.167) (6.5,2.000) (7.5,1.875) (8.5,2.667) (9.5,2.500) (10.5,2.364) (11.5,2.250) (12.5,2.154) (13.5,2.071)
    (14.5,2.000) (15.5,1.938) (16.5,2.824) (17.5,2.722) (18.5,2.632) (19.5,2.550) (20.5,2.476) (21.5,2.409) (22.5,2.348)
    (23.5,2.292) (24.5,2.240) (25.5,2.192) (26.5,2.148) (27.5,2.107) (28.5,2.069) (29.5,2.033) (30.5,2.000) (31.5,1.969)};
  \draw[line width=0.8pt, fgink, dashed] (0, 3) -- (32, 3) node[right, font=\sffamily\scriptsize] {3};
  \draw[->, line width=0.6pt, fgink!70] (0, 0) -- (33, 0) node[right, lbl] {push number};
  \draw[->, line width=0.6pt, fgink!70] (0, 0) -- (0, 19) node[above, lbl] {cost};
  \foreach \k in {1, 8, 16, 24, 32} \node[lbl, below] at (\k-0.5, 0) {\k};
  \foreach \v in {5, 10, 15} \node[lbl, left] at (0, \v) {\v};
  \node[font=\sffamily\scriptsize, text=fgverm, anchor=south] at (16.5, 17.2) {copy 16};
  \node[font=\sffamily\scriptsize, text=fggreen!80!black, anchor=west] at (20, 5.2) {average cost so far};
\end{tikzpicture}
\caption{The cost of 32 \texttt{push\_back} calls with capacity doubling. A few pushes are expensive (red: they copy the
whole vector), but the average cost per push (green) never exceeds 3.}
\label{fig:doubling}
\end{figure}

The same argument explains the two-pointer loop below. The inner \cc{while} can run many times for a single \cc{r}, but
\cc{l} only ever moves forward and never passes $n$, so over the whole loop the inner body runs at most $n$ times in total.
The loop is $\Oh{n}$, not $\Oh{n^2}$.

\begin{code}
for (int l = 0, r = 0; r < n; r++) {
    add(a[r]);
    while (bad()) remove(a[l++]);  // l moves at most n times in total
}
\end{code}

\subsection{The potential method}

For more complicated structures, define a \defn{potential} $\Phi \ge 0$ that measures ``stored-up work'', and charge each
operation its real cost plus the change in potential: $\hat c = c + \Phi_{\text{after}} - \Phi_{\text{before}}$. Summed
over a sequence, the potentials telescope, so the total real cost is at most the total charged cost (as long as $\Phi$
starts at 0 and never goes negative). The art is choosing $\Phi$ so that every operation's charged cost is small.

For a monotonic stack (Chapter~\ref{ch:prefix}), take $\Phi$ = the number of elements on the stack. A push costs 1 and
raises $\Phi$ by 1: charged cost 2. Popping $k$ elements costs $k$ and lowers $\Phi$ by $k$: charged cost 0. Every
operation is charged at most 2, so $n$ operations cost at most $2n$ in total, however the pops are distributed. The same
method proves the $\Oh{\log n}$ amortized bounds of splay trees, the linear time of the KMP prefix function, and the nearly
constant time of union--find with path compression.

\begin{pitfall}[Amortized is not worst case]
An amortized bound holds for the whole sequence, not for each operation. If a problem has multiple test cases and you
``reset'' a structure by rebuilding it, or if a data structure is persistent and an old expensive state can be revisited
many times, the amortization can break. Rollback DSU (Chapter~\ref{ch:dsu}) avoids path compression for exactly this reason.
\end{pitfall}

% ======================================================================================================================
\section[Recursion]{Recursion: trees and the master theorem \hfill \gold}\label{sec:recursion}

A recursive algorithm that splits a problem of size $n$ into $a$ subproblems of size $n/b$ and does $f(n)$ extra work
(splitting and combining) takes time
\[ T(n) = a\,T(n / b) + f(n). \]
The reliable way to solve such a recurrence is to draw its \defn{recursion tree}: the root is the original call, each node
has $a$ children, and each node is labelled with its own extra work. Adding up the work level by level gives the total
(Figure~\ref{fig:mergesort-tree}).

\begin{figure}[htbp]
\centering
\begin{tikzpicture}[lvl/.style={segn, minimum height=5.5mm}, lab/.style={font=\sffamily\scriptsize, text=fgink, anchor=west}]
  \node[lvl, minimum width=11.4cm] at (0, 0) {$n$};
  \foreach \i in {0,1} \node[lvl, minimum width=5.6cm] at (-2.9 + \i*5.8, -0.85) {$n/2$};
  \foreach \i in {0,...,3} \node[lvl, minimum width=2.7cm] at (-4.35 + \i*2.9, -1.7) {$n/4$};
  \foreach \i in {0,...,7} \node[lvl, minimum width=1.25cm] at (-5.075 + \i*1.45, -2.55) {$n/8$};
  \node[font=\small, text=fggray] at (0, -3.2) {$\vdots$};
  \foreach \i in {0,...,15} \node[seg, minimum width=0.6cm, minimum height=5mm] at (-5.4375 + \i*0.725, -3.85) {\tiny 1};
  \foreach \y/\t in {0/{$n$}, -0.85/{$2 \cdot \frac n2 = n$}, -1.7/{$4 \cdot \frac n4 = n$}, -2.55/{$8 \cdot \frac n8 = n$},
      -3.85/{$n \cdot 1 = n$}}
    \node[lab] at (6.0, \y) {\t};
  \node[font=\sffamily\scriptsize\bfseries, text=fgink, anchor=west] at (6.0, 0.6) {work per level};
  \draw[brace] (-6.1, -4.1) -- (-6.1, 0.3) node[midway, left=8pt, font=\sffamily\scriptsize, align=center] {$\log_2 n$\\levels};
\end{tikzpicture}
\caption{The recursion tree of merge sort, $T(n) = 2T(n/2) + \Theta(n)$. Every level does $\Theta(n)$ work in total and
there are $\log_2 n$ levels, so $T(n) = \Theta(n \log n)$.}
\label{fig:mergesort-tree}
\end{figure}

Whether the root, the leaves, or all levels together dominate depends on how the work per level changes as you go down
(Figure~\ref{fig:master}). That gives the three cases of the \defn{master theorem}:

\begin{fact}[Master theorem]
Let $T(n) = a\,T(n/b) + \Theta(n^c)$ and compare $c$ with the critical exponent $\log_b a$.
\begin{itemize}
  \item If $c > \log_b a$, the work shrinks down the tree and the root dominates: $T(n) = \Theta(n^c)$.
  \item If $c = \log_b a$, every level does the same work: $T(n) = \Theta(n^c \log n)$.
  \item If $c < \log_b a$, the work grows down the tree and the leaves dominate: $T(n) = \Theta(n^{\log_b a})$.
\end{itemize}
\end{fact}

\begin{figure}[htbp]
\centering
\begin{tikzpicture}[lvbar/.style={line width=0.4pt, draw=#1, fill=#1!35}]
  \foreach \panel/\x/\title/\res/\col/\ws in {
      1/0/{$T(n) = T(n/2) + n$}/{root dominates: $\Theta(n)$}/fgblue/{8,4,2,1,0.5},
      2/4.9/{$T(n) = 2T(n/2) + n$}/{all equal: $\Theta(n \log n)$}/fggreen/{8,8,8,8,8},
      3/9.8/{$T(n) = 4T(n/2) + n$}/{leaves dominate: $\Theta(n^2)$}/fgverm/{1,2,4,8,16}} {
    \begin{scope}[xshift=\x cm]
      \node[font=\sffamily\scriptsize\bfseries] at (1.8, 0.55) {\title};
      \foreach \w [count=\lv from 0] in \ws {
        \pgfmathsetmacro{\len}{\w*0.22}
        \fill[fill=\col!35, draw=\col, line width=0.4pt] (0, -\lv*0.5) rectangle (\len, -\lv*0.5-0.34);
        \node[lbl, anchor=east] at (-0.05, -\lv*0.5-0.17) {\lv};
      }
      \node[font=\sffamily\scriptsize, text=\col!80!black] at (1.8, -2.75) {\res};
    \end{scope}
  }
  \node[lbl, anchor=east] at (-0.4, 0.1) {level};
  \node[lbl] at (6.7, -3.25) {bar length = total work on that level};
\end{tikzpicture}
\caption{Work per level of the recursion tree in the three cases of the master theorem. Geometric decrease means the top
level dominates; a constant means $\log n$ equal levels; geometric increase means the leaves dominate.}
\label{fig:master}
\end{figure}

Some recurrences do not fit the theorem but are easy with a tree anyway. Linear recursion $T(n) = T(n-1) + \Theta(n)$ (a
quicksort that always picks the worst pivot, or selection sort) has $n$ levels with $n, n-1, \dots$ work: $\Theta(n^2)$.
$T(n) = 2T(n-1) + 1$ doubles the number of calls per level: $\Theta(2^n)$, which is why naive recursive Fibonacci is
hopeless. Memoised recursion is analysed differently: its cost is (number of distinct states) $\times$ (work per state),
the same as the equivalent DP (Chapter~\ref{ch:dp}).

\begin{center}
\small
\begin{tabularx}{\linewidth}{Lll}
\toprule
\thd{Algorithm} & \thd{Recurrence} & \thd{Solution} \\
\midrule
binary search & $T(n) = T(n/2) + \Theta(1)$ & $\Theta(\log n)$ \\
merge sort, divide-and-conquer DP & $T(n) = 2T(n/2) + \Theta(n)$ & $\Theta(n \log n)$ \\
building a segment tree bottom-up & $T(n) = 2T(n/2) + \Theta(1)$ & $\Theta(n)$ \\
Karatsuba multiplication & $T(n) = 3T(n/2) + \Theta(n)$ & $\Theta(n^{\log_2 3}) \approx \Theta(n^{1.585})$ \\
Strassen matrix multiplication & $T(n) = 7T(n/2) + \Theta(n^2)$ & $\Theta(n^{\log_2 7}) \approx \Theta(n^{2.81})$ \\
CDQ divide and conquer with a Fenwick tree & $T(n) = 2T(n/2) + \Theta(n \log n)$ & $\Theta(n \log^2 n)$ \\
naive recursive Fibonacci & $T(n) = T(n-1) + T(n-2) + \Theta(1)$ & $\Theta(\varphi^n)$, $\varphi \approx 1.618$ \\
\bottomrule
\end{tabularx}
\end{center}

% ======================================================================================================================
\section[Memory]{Memory \hfill \bronze}\label{sec:memory}

Memory limits are checked as strictly as time limits, and running out gives \emph{memory limit exceeded} or a runtime error.
Budgeting memory is simple multiplication once you know the sizes.

\begin{center}
\small
\begin{tabularx}{\linewidth}{lrL}
\toprule
\thd{Type} & \thd{Bytes} & \thd{How many fit in 256\,MB} \\
\midrule
\texttt{char}, \texttt{bool} & 1 & $2.6\cdot 10^8$ (a \texttt{vector<bool>} or \texttt{bitset} stores 8 per byte) \\
\texttt{int32\_t}, \texttt{float} & 4 & $6.7\cdot 10^7$ \\
\texttt{long long}, \texttt{double} & 8 & $3.3\cdot 10^7$ \\
\texttt{pair<ll, ll>}, \texttt{array<ll, 2>} & 16 & $1.6\cdot 10^7$ \\
\texttt{vector} header (empty vector) & 24 & plus the heap block it owns \\
\texttt{set} / \texttt{map} node & about 40--48 & about $5\cdot 10^6$ elements \\
\texttt{unordered\_map} entry & about 40--64 & fewer than a \texttt{map}; plus the bucket array \\
\bottomrule
\end{tabularx}
\end{center}

\begin{pitfall}[Your template doubles every int]
With \cc{#define int long long}, a \cc{vector<int>} of $10^7$ elements takes 80\,MB instead of 40\,MB, and a DP table of
$5000 \times 5000$ ints takes 200\,MB. For big arrays that only hold small values, write \cc{int32_t} or \cc{char}
explicitly; the macro does not touch them.
\end{pitfall}

Two memory traps have nothing to do with element sizes. The first is \emph{many small vectors}: a
\cc{vector<vector<int>>} with $10^6$ rows spends 24 bytes per row on headers alone, plus a separate heap allocation (with
its own overhead of 16 bytes or more) for every non-empty row, and its rows are scattered across memory, which also makes
it slow (Figure~\ref{fig:vecvec}). When the rows are fixed after reading the input, a flat array with offsets (compressed
sparse rows) uses a fraction of the memory and is several times faster to scan.

\begin{figure}[htbp]
\centering
\begin{tikzpicture}[hdr/.style={draw=fgblue, fill=fgfill, line width=0.6pt, minimum width=1.9cm, minimum height=5.5mm,
    font=\sffamily\scriptsize, inner sep=1pt}, blk/.style={draw=fgverm, fill=fgverm!15, line width=0.6pt, minimum height=5.5mm,
    font=\ttfamily\scriptsize, inner sep=1pt}]
  \node[font=\small\bfseries, anchor=west] at (-1.0, 0.7) {\texttt{vector<vector<int>{}>}};
  \foreach \i in {0,1,2} \node[hdr] (h\i) at (0, -\i*0.62) {ptr, size, cap};
  \node[blk, minimum width=2.1cm] (b0) at (4.6, 0.3) {3 1 4};
  \node[blk, minimum width=1.4cm] (b1) at (6.2, -0.62) {1 5};
  \node[blk, minimum width=2.8cm] (b2) at (3.9, -1.55) {9 2 6 5};
  \draw[->, line width=0.7pt, fgink!70] (h0.east) -- (b0.west);
  \draw[->, line width=0.7pt, fgink!70] (h1.east) -- (b1.west);
  \draw[->, line width=0.7pt, fgink!70] (h2.east) -- (b2.west);
  \node[lbl, text width=4.8cm, align=left, anchor=north west] at (-1.0, -2.25) {24-byte header per row, one heap block per
    row (scattered), allocator overhead on each block};
  \begin{scope}[xshift=9.0cm]
    \node[font=\sffamily\small\bfseries, anchor=west] at (-0.5, 0.7) {flat array + offsets};
    \arr{F}{-0.5}{0}{3,1,4,1,5,9,2,6,5}
    \arr{O}{-0.5}{-1.15}{0,3,5,9}
    \node[lbl, anchor=east] at (-0.6, 0) {data};
    \node[lbl, anchor=east] at (-0.6, -1.15) {start};
    \node[lbl, text width=5.4cm, align=left, anchor=north west] at (-0.5, -1.85) {row $i$ is
      \texttt{data[start[i] .. start[i+1])}: one allocation, contiguous, cache friendly};
  \end{scope}
\end{tikzpicture}
\caption{The same three rows stored two ways. A vector of vectors pays a header and an allocation per row and scatters the
rows; a flat array with offsets stores only the data and one integer per row.}
\label{fig:vecvec}
\end{figure}

The second trap is the \emph{stack}. Local arrays and recursion live on the call stack, which on many judges is limited to
8\,MB or less, though Codeforces and your local setup use 256\,MB. A recursive DFS on a path of $10^6$ vertices needs a
million stack frames; with 50--100 bytes per frame it fits in 256\,MB but not in 8\,MB. Big arrays belong in global
variables or vectors, never in local arrays; deep recursion should become an explicit stack when the judge is unknown.

% ======================================================================================================================
\section[How fast is an operation, really?]{How fast is an operation, really? \hfill \gold}\label{sec:measured}

``$10^8$ operations per second'' is an average over very different operations. Table~\ref{tab:measured} was measured on
a 2026 laptop with GCC 15 and \texttt{-O2}, the same flags as most judges; judges differ from it by a factor of about two in
either direction, but the \emph{ratios} are what matter, and they are stable across machines.

\begin{table}[htbp]
\centering
\caption{Measured cost of common operations (GCC 15, \texttt{-O2}, averaged over $10^6$ to $10^9$ repetitions).}
\label{tab:measured}
\small
\begin{tabularx}{\linewidth}{Lrr}
\toprule
\thd{Operation} & \thd{Time per operation} & \thd{Operations per second} \\
\midrule
reading an array sequentially (the compiler vectorizes the loop) & 0.1 ns & $10^{10}$ \\
dependent add / xor / shift in a loop & 0.6 ns & $1.6 \cdot 10^9$ \\
\texttt{vector::push\_back} (amortized) & 1.5 ns & $6.6 \cdot 10^8$ \\
\texttt{x = x * c \% MOD} with 64-bit numbers & 2.9 ns & $3.4 \cdot 10^8$ \\
\texttt{std::sort} of $10^6$ ints, per element & 43 ns & $2.3 \cdot 10^7$ \\
\texttt{priority\_queue} push + pop, per element & 60 ns & $1.7 \cdot 10^7$ \\
reading one integer with \texttt{cin} (sync off) & 80 ns & $1.2 \cdot 10^7$ \\
random access into a 64\,MB array (cache miss) & 88 ns & $1.1 \cdot 10^7$ \\
\texttt{unordered\_map::operator[]}, with / without \texttt{reserve} & 135 / 215 ns & $5$--$7 \cdot 10^6$ \\
\texttt{set::insert}, \texttt{map::operator[]} & 400--440 ns & $2.4 \cdot 10^6$ \\
\texttt{set::count} on a set of $10^6$ random keys & 470--500 ns & $2 \cdot 10^6$ \\
\bottomrule
\end{tabularx}
\end{table}

Four lessons stand out. \textbf{Memory access patterns dominate}: a random access into a large array costs as much as a
hundred additions, because it misses the cache and waits for main memory. Scanning arrays in order is nearly free.
\textbf{Balanced trees are slow}: every \cc{set} or \cc{map} operation is a chain of dependent cache misses down a tree, so
$10^6$ operations take almost half a second; a sorted \cc{vector} with \cc{lower_bound}, or an array indexed by compressed
values, is often ten times faster. \textbf{Division is expensive} relative to addition: reduce modulo only when needed.
\textbf{I/O is not free}: reading $10^6$ numbers takes 0.08\,s even with fast settings, and many times more without them.

% ======================================================================================================================
\section[Constant factors that matter]{Constant factors that matter \hfill \gold}\label{sec:constants}

Two $\Oh{n \log n}$ solutions can differ by a factor of twenty in running time, and near the limit that factor decides
between accepted and \emph{time limit exceeded}. The usual culprits, in order of how often they matter:

\begin{itemize}
  \item \textbf{Slow I/O.} Without \cc{ios::sync_with_stdio(false)} and \cc{cin.tie(nullptr)}, reading $10^6$ numbers can
        take a second; \cc{endl} flushes the output buffer on every line. Your template sets up the fast mode; print
        \cc|'\n'|, never \cc{endl}, in loops.
  \item \textbf{Node-based containers.} Replace \cc{set}/\cc{map} by sorted vectors, prefix sums or arrays over
        compressed values whenever the data does not change after it is built.
  \item \textbf{Cache-hostile loops.} Traverse two-dimensional arrays row by row, never column by column. In matrix
        multiplication, the loop order $i, k, j$ is several times faster than $i, j, k$.
  \item \textbf{Needless modulo.} Add, then subtract \cc{MOD} if the sum reached it; accumulate several products in a
        64-bit variable and reduce once.
  \item \textbf{Copies.} Passing a \cc{vector} or \cc{string} by value copies it on every call; pass \cc{const&}.
  \item \textbf{Hidden allocation.} Creating a \cc{vector} inside a hot loop allocates every time; hoist it out and
        \cc{clear()} it.
\end{itemize}

\begin{pitfall}[Hidden linear costs]
These look like single operations but cost $\Oh{n}$ each: \cc{s = s + c} on a string (use \cc{s += c}),
\cc{v.erase(v.begin())} and \cc{v.insert(v.begin(), x)}, \cc{std::distance} or \cc{std::next} on a \cc{set} iterator,
\cc{count} or \cc{find} on a \cc{vector}, \cc{memset} of a whole array in every test case, \cc{s.substr(...)} inside a loop,
and \cc{mp[key]}, which \emph{inserts} a default element when the key is missing and slowly fills the map.
\end{pitfall}

\begin{trick}[Bitsets divide by 64]
A \cc{bitset<N>} stores bits in 64-bit words, and \cc{&}, \cc{|}, \cc{^}, shifts and \cc{count()} process 64 bits per
instruction. Any computation that is ``OR of shifted boolean arrays'' becomes 64 times faster: subset sum up to $10^5$ over
$10^4$ items needs $10^9/64 \approx 1.6\cdot10^7$ word operations. Reachability in a DAG, triangle detection and
longest common subsequence have bitset versions too (Chapter~\ref{ch:special}).
\end{trick}

% ======================================================================================================================
\section[Estimating before coding]{Estimating before coding: a worked example \hfill \silver}\label{sec:estimate}

Suppose a problem has $n = 2 \cdot 10^5$ elements and $q = 2 \cdot 10^5$ queries, and you have three ideas.

\begin{enumerate}
  \item \textbf{Answer each query by scanning the array}: $nq = 4 \cdot 10^{10}$. Rejected immediately, without coding.
  \item \textbf{Mo's algorithm with a \cc{map} inside}: $(n + q)\sqrt n$ pointer moves, each doing a \cc{map} operation:
        $4\cdot10^5 \cdot 450 \approx 1.8\cdot 10^8$ moves at 400\,ns each is 72 seconds. Rejected: the complexity was
        fine but the constant is not (Section~\ref{sec:measured}). With an array instead of the \cc{map}, each move costs
        a few nanoseconds and the same idea runs in under a second.
  \item \textbf{Offline sweep with a Fenwick tree}: $(n + q)\log n \approx 4\cdot10^5 \cdot 18 = 7\cdot10^6$ simple
        operations. Clearly fast; code this one.
\end{enumerate}

Doing this arithmetic takes a minute and saves you from implementing the wrong idea. Notice that idea 2 was rejected for
its constant, not its class: once you know the measured cost of your operations, ``fast enough'' is a calculation.

\begin{summary}[Chapter summary]
\begin{itemize}
  \item Budget about $10^8$ simple operations per second; count the operations of your idea with the real limits.
  \item Read the intended complexity from the limits (Table~\ref{tab:constraints}); the smallest limit is often the hint.
  \item Harmonic sums, halving loops, submask loops and amortized arguments are cheaper than they look.
  \item Solve recurrences with a recursion tree; the master theorem compares $c$ with $\log_b a$.
  \item Memory is multiplication: know your type sizes, beware vectors of vectors and the stack.
  \item Constants matter near the limit: cache misses, tree containers, modulo and I/O cost far more than an addition.
\end{itemize}
\end{summary}

\begin{practice}
\ppgroup{Estimate, then code}
\begin{pplist}
  \item Weird Algorithm
  \item Distinct Numbers
  \item Sum of Two Values
  \item Josephus Problem I
  \item Concert Tickets
  \item Reading Books
\end{pplist}
\ppgroup{Harmonic sums and square roots}
\begin{pplist}
  \item Counting Divisors
  \item Common Divisors
  \item Sum of Divisors
  \item Prime Multiples
\end{pplist}
\end{practice}
